\documentclass[12pt]{article}
\usepackage{amsmath,amssymb}
\newcommand{\conv}{\operatorname{conv}}
\newcommand{\Ber}{\operatorname{Bernoulli}}
\begin{document}
\section*{Research notes}

\subsection*{Problem interpretation and expected answer}

Universal interpretation: determine all \(p\in[0,1]\) for which
\[
 \Pr\left[\sum_i w_i v_i\ge p\right]\ge p
\]
holds for every finite nonnegative probability vector \(w\), with independent
\(v_i\sim{\rm Bernoulli}(p)\).

Current expected answer:
\[
 [0,1/3]\cup\{1/2,1\}.
\]
The intervals outside this set are excluded by equal-weight examples:
\begin{itemize}
\item \(m=3\), \(w_i=1/3\), excludes \(1/3<p<1/2\), since the event needs at
  least two successes and has probability \(3p^2-2p^3<p\).
\item \(m=2\), \(w_i=1/2\), excludes \(1/2<p<1\), since the event needs two
  successes and has probability \(p^2<p\).
\end{itemize}
The points \(p=0,1/2,1\) are elementary.  At \(p=1/2\), \(X\stackrel d=1-X\).

\subsection*{Reciprocal points}

For \(p=1/k\), \(k\ge3\), a complete coloring proof works.  Color each index
uniformly in \(k\) colors.  Let \(W_j\) be the total weight of color \(j\).  Since
\(\sum_j W_j=1\), some \(W_j\ge1/k\).  By symmetry,
\(\Pr[W_1\ge1/k]\ge1/k\), and color \(1\) is an independent Bernoulli\((1/k)\)
selection.  This proves all reciprocal points, in particular \(p=1/3\).

\subsection*{Main missing lemma}

Need prove the finite weighted Bernoulli inequality:
for \(0\le p\le1/3\), \(w_i\ge0\), \(\sum_i w_i=1\),
\[
 \Pr\left[\sum_i w_i v_i\ge p\right]\ge p.
\]
If some \(w_i\ge p\), the event contains \(\{v_i=1\}\), so the result is
immediate.  Thus the nontrivial case has all \(w_i<p\).  Normalize
\[
 a_i=w_i/p,\qquad L=1/p\ge3.
\]
Then \(0\le a_i<1\), \(\sum_i a_i=L\), and \(\xi_i\sim{\rm Bernoulli}(1/L)\);
the target becomes
\[
 \Pr\left[\sum_i a_i\xi_i\ge1\right]\ge\frac1L .
\]
For integral \(L\) this is the reciprocal coloring argument.  Non-integral
\(L>3\) is the remaining issue.

Complement form: with \(q=1-p\ge2/3\), the theorem is equivalent to
\[
 \Pr\left[\sum_i w_i \eta_i>q\right]\le q,\qquad
 \eta_i\sim{\rm Bernoulli}(q),
\]
because \(\sum_i w_i v_i\ge p\) iff \(\sum_i w_i(1-v_i)\le1-p=q\).
If some \(w_i>q\), the complement event is exactly \(\{\eta_i=1\}\), since
the remaining total weight is \(<1-q\le q\); hence probability \(q\).

\subsection*{Attempted reductions and checks}

\paragraph{Rational-color deterministic strengthening fails.}
For \(p=a/b\), one can generate Bernoulli\((a/b)\) variables by coloring each
coordinate uniformly in \(b\) colors and selecting \(a\) colors.  A tempting
deterministic claim would be: for every nonnegative \(b\)-tuple \(y_1,\ldots,y_b\)
with sum \(1\), at least an \(a/b\) fraction of all \(a\)-subsets have sum at
least \(a/b\).  This is false even with \(a/b\le1/3\): take \(a=3,b=10\),
seven entries \(1/7\) and three entries \(0\).  A 3-subset is good iff it contains
at least three of the seven positive entries, so the fraction is
\[
 \binom73/\binom{10}3=35/120<3/10.
\]
Thus a rational-color proof must exploit randomness of the color weights induced
by the original independent coloring, not an arbitrary deterministic color vector.

\paragraph{Circle interval deterministic strengthening fails.}
Another tempting representation is \(v_i=\mathbf1_{\{U_i\in[t,t+p]\}}\) on the
circle and to try to prove that for every atomic probability measure on the
circle, the set of interval starts \(t\) for which the length-\(p\) arc has
mass at least \(p\) has Lebesgue measure at least \(p\).  Numerical exact checks
show this deterministic statement is false.  Example for \(p=0.3\): atoms near
\(0.2093,0.5049,0.8102,0.9471\) with weights approximately
\(0.2755,0.1601,0.2716,0.2928\) have good-start measure about \(0.2053<0.3\).
Therefore any circle proof must use averaging over the random atom positions in
an essential way.

\paragraph{Splitting/merging is not monotone pointwise.}
Splitting one Bernoulli weight \(c=a+b\) into independent pieces \(a,b\) can
decrease or increase the probability of crossing a fixed threshold, depending on
the rest of the configuration.  Conditional on a rest value that leaves residual
threshold \(r\), old contribution \(c\zeta\) succeeds with probability \(p\) for
\(0<r\le c\); split contribution \(a\xi+b\eta\) has probability \(p^2\), \(p\),
or \(2p-p^2\) according to the location of \(r\).  Example at \(p=0.3\):
\[
 \Pr_{[0.7,0.2,0.1]}(X\ge0.3)=0.363,\qquad
 \Pr_{[0.7,0.1,0.1,0.1]}(X\ge0.3)=0.3189.
\]
Both are above \(p\), but simple monotonicity under splitting is false.

\paragraph{Moment and polynomial lower bounds tried.}
Pointwise polynomial minorants of \(\mathbf1_{\{x\ge p\}}\), e.g.
\(x(x-p)/(1-p)\), give
\[
 \Pr[X\ge p]\ge \frac{\mathbb E[X^2]-p^2}{1-p}=p\sum_i w_i^2,
\]
which is sharp for a single atom but too weak for spread-out weights.  Cantelli
or Bhatia--Davis style variance arguments similarly give only \(p\sum_i w_i^2\).
Chernoff bounds at the mean are trivial.  Concave-transform minorants such as
\(1-((1-x)/(1-p))^\alpha\) give nonnegative bounds but do not reach \(p\) in the
single-atom equality case except in a limiting sense.

\subsection*{Possible future proof directions}

\begin{itemize}
\item Find the correct Kellerer theorem/reference.  The previously cited paper
  ``Zur Existenz analoger Bereiche'' proves existence of analogous regions for
  finite atomless measure families and does not directly state the Bernoulli
  halfspace inequality.  If it is relevant, a derivation must specify the
  atomless embedding and how the weak threshold is preserved.
\item Look for a theorem on sums \(S=\sum_i a_i\xi_i\) with \(0\le a_i\le1\),
  \(\sum a_i=L\ge3\), \(\xi_i\sim{\rm Bernoulli}(1/L)\), asserting
  \(\Pr[S\ge1]\ge1/L\).  Search terms: weighted Bernoulli mean-tail inequality,
  probability above the mean for weighted Bernoulli sums, Kellerer inequality,
  Feige/Kleitman small deviation inequalities, p-biased measure of weighted
  halfspaces.
\item Try a minimal-counterexample proof.  The objective is piecewise polynomial
  in weights, with changes only when subset sums cross \(p\).  At a minimizer
  one may be able to move in the simplex until a boundary \(w_i=0\), \(w_i=p\),
  or a subset-sum equality is reached; induction over such faces may be possible.
\item Try an induction for the stronger two-parameter tail function
  \(H_{p}(s,t)=\inf\Pr[\sum w_i v_i\ge t]\) over weights of total \(s\).  The
  one-parameter induction at \(s=1,t=p\) fails because conditioning on a weight
  produces thresholds \(p/(1-a)>p\) and \((p-a)/(1-a)<p\).
\item Search for a coupling or injection proving
  \(\Pr(\sum_i w_i v_i\ge p)\ge \Pr(I\in R)=p\), where \(I\) is an independent
  size-biased index with \(\Pr(I=i)=w_i\) and \(R\) is the \(p\)-biased random
  subset.  Direct containment \(\{I\in R\}\subseteq\{w(R)\ge p\}\) is false, but
  a measure-preserving repair map might exist for \(p\le1/3\).
\end{itemize}

\subsection*{References noted}

H. G. Kellerer, ``Zur Existenz analoger Bereiche'', Z. Wahrscheinlichkeitstheorie
verw. Gebiete 1 (1963), 240--246.  This reference should not be used as a black
box for the Bernoulli inequality without an explicit derivation.

Related but not directly sufficient: Manickam--Miklós--Singhi nonnegative
\(k\)-sums and Baranyai partition arguments.  Divisible cases match the
reciprocal coloring proof; non-divisible deterministic sampling statements can
fail and hence do not settle the Bernoulli problem.


\subsection*{Round 2 additional structural reductions}

\paragraph{Finite convex-hull/down-set equivalence.}
For fixed \(p\), the small-\(p\) lemma is equivalent to the following finite
``Kellerer-type'' convex-hull statement: if \(\mathcal D\subseteq 2^{[m]}\) is a
down-set and \(\mu_p(\mathcal D)>1-p\), then
\[
 (p,\ldots,p)\in\operatorname{conv}\{\mathbf 1_A:A\in\mathcal D\}.
\]
Implication to Bernoulli inequality: for weights \(w\), set
\(\mathcal D=\{A:w(A)<p\}\).  If \(\mu_p(\mathcal D)>1-p\), a distribution
\(\nu\) on \(\mathcal D\) with coordinate marginals \(p\) gives
\[
p=\sum_i p w_i=\mathbb E_\nu w(A)<p.
\]
Conversely, if the convex-hull statement fails, since the convex hull of a
down-set is itself down-closed, separating it from
\((p,\ldots,p)+\mathbb R_{\ge0}^m\) gives a nonnegative vector \(c\) such that
\(c(A)<p\sum_i c_i\) for all \(A\in\mathcal D\).  Normalize
\(w_i=c_i/\sum_j c_j\).  The Bernoulli inequality implies
\(\mu_p(\mathcal D)\le \Pr[w(R)<p]\le1-p\).

This equivalence is useful but also warns against proving the convex-hull lemma
only by separation, since the separated inequality is exactly the weighted
Bernoulli inequality.

\paragraph{Complement/Fokkink--Meester--Pelekis formulation.}
With \(q=1-p\), the desired \(p\le1/3\) statement is equivalent to
\[
 \Pr\left[\sum_i w_i\eta_i>q\right]\le q,
 \qquad q\in[2/3,1],\quad \eta_i\sim{\rm Bernoulli}(q).
\]
In the notation \(\pi(q,t)=\sup_\gamma\Pr(S_\gamma\ge t)\) for convex
combinations of iid Bernoulli\((q)\)'s, it would follow from
\[
 \pi(q,t)\le q\qquad(q\ge2/3,\ t>q).
\]
Fokkink--Meester--Pelekis prove adjacent ``bold play'' regions but their
published theorems leave near-hypotenuse triangles; their results should not be
cited as proving this full region without an additional argument.

\paragraph{Weighted quantile formulation.}
Let \(U_i\) be iid uniform on \([0,1]\), and let
\[
 Q_p=\inf\{t:\sum_{i:U_i\le t} w_i\ge p\}.
\]
Then
\[
 \Pr\left[\sum_i w_i v_i\ge p\right]=\Pr[Q_p\le p].
\]
After ordering the \(U_i\)'s, the weights appear in a uniformly random
permutation.  If \(K\) is the first position in this random permutation whose
cumulative weight reaches \(p\), then
\[
 \Pr[Q_p\le p]=\mathbb E\,\Pr[{\rm Bin}(m,p)\ge K\mid K].
\]
For a single weight \(1\) (and any number of zero weights), \(K\) is uniform and
the expectation is exactly \(p\).  For equal small weights, \(K\) is near \(pm\)
and the probability is near \(1/2\).  This supports the heuristic that the
minimum is approached by a nearly single atom of weight \(p\), but a proof needs
a stochastic order or extremal theorem for this weighted hitting position.

\subsection*{New positive family \(p=2/b\)}

For \(b\ge6\), the deterministic pair lemma holds: if \(x_1,\ldots,x_b\ge0\)
and \(\sum x_i=1\), then at least \(b-1\) pairs have
\(x_i+x_j\ge2/b\).  Scaling \(z_i=bx_i/2\), this becomes: if
\(\sum z_i=b/2\), then at least \(b-1\) pairs have \(z_i+z_j\ge1\).

Proof sketch retained for future use.  Given a remaining set \(R\) of size
\(s\) and sum at least \(s/2\), let \(u\) and \(a\) be its minimum and maximum.
If \(u+a\ge1\), the maximum forms \(s-1\) good pairs and one can stop.  If
\(u+a<1\), let \(h\) be the number of partners of \(a\).  If
\(h\le(s-2)/2\), then
\[
 \sum_{z\in R}z
 < a+ha+(s-1-h)(1-a)
 =s-1-h+a(2h-s+2)\le s/2,
\]
contradiction.  Thus \(h\ge\lfloor s/2\rfloor\).  Count those pairs, delete
\(u,a\), and iterate; the residual sum remains \(>(s-2)/2\).  Summing the
contributions gives at least \(b-1\) for \(b\ge6\).

Consequently \(p=2/b\) is positive: color each original index uniformly in
\(b\) colors; by the pair lemma every color-load vector has at least \(b-1\)
good color-pairs.  Symmetry gives
\[
 \Pr[\text{colors }1\text{ or }2\text{ have load}\ge2/b]
 \ge \frac{b-1}{\binom b2}=2/b.
\]
Membership in colors \(1\) or \(2\) is independent Bernoulli\((2/b)\).

The pair lemma fails for \(b=5\), e.g. after scaling back a vector close to
\((0.39,0.39,0.21,0.005,0.005)\) has only three good pairs while \(b-1=4\).
This matches the global threshold: \(2/5>1/3\) is outside the positive
interval.

\subsection*{Further computational/heuristic observations}

\paragraph{Random quantile and splitting.}
Splitting a weight can decrease the success probability even below the merged
configuration, e.g. for \(p=0.3\) splitting one of two large weights can move
probability from \(0.51\) to \(0.363\).  However, in all numerical searches for
\(p\le1/3\) the value remained above \(p\), with infima approached by one weight
near \(p\) and a finely split residual.

\paragraph{Influence/differential route.}
For a fixed threshold family \(F_t=\{A:w(A)\ge t\}\), \(\mu_r(F_t)\) has
derivative equal to total influence.  The moving-threshold function
\(f(p)=\mu_p(\{w(A)\ge p\})-p\) is polynomial between subset sums and has
downward jumps just after subset-sum values.  A possible proof would need both
a lower bound on the total influence between jumps and a control of the jump
sizes at \(w(A)=p\).  The total influence bound is false for arbitrary up-sets,
but may hold for these near-mean weighted threshold up-sets.


\subsection*{Multiplicative closure}

Let \(\mathcal G\) denote the set of probabilities for which the universal
inequality is true.  A useful closure property is
\[
        p,q\in\mathcal G \quad\Longrightarrow\quad pq\in\mathcal G .
\]
Proof: generate Bernoulli\((pq)\) variables as products \(x_i y_i\), with
\(x_i\sim{\rm Bernoulli}(p)\) and \(y_i\sim{\rm Bernoulli}(q)\) independent.
With \(T=\sum_i w_i x_i\), the event \(T\ge p\) has probability at least \(p\).
Conditional on such an outcome and normalizing the selected weights by \(T\),
the \(q\)-property gives conditional probability at least \(q\) that the thinned
sum is at least \(qT\ge pq\).  Hence the product probability works.

Consequences:
\begin{itemize}
\item Products of any already positive values are positive.
\item If one proves the small-\(p\) lemma on the base interval
  \([1/6,1/3]\), then all of \([0,1/3]\) follows by multiplying by powers of
  \(1/2\).  Indeed, for any \(0<p\le1/3\), some \(2^n p\in[1/6,1/3]\).
\item More generally, if a base set \(B\subseteq\mathcal G\) multiplicatively
  covers \((0,1/3]\) after products with known values, it suffices to prove
  \(B\).
\end{itemize}

\subsection*{Partition-function target for the convex-hull form}

For a down-set \(\mathcal D\subseteq2^{[m]}\), put
\[
        Z_\lambda(\mathcal D)=\sum_{A\in\mathcal D}\lambda^{|A|},
        \qquad \lambda={p\over 1-p}.
\]
The convex-hull formulation of the missing lemma is equivalent to
\[
 p{\bf1}\notin \operatorname{conv}\{\mathbf1_A:A\in\mathcal D\}
 \quad\Longrightarrow\quad
 Z_\lambda(\mathcal D)\le (1+\lambda)^{m-1}
\]
for \(0\le\lambda\le1/2\).  The extremal example is a dictatorship
\(\mathcal D=\{A:i\notin A\}\), where equality holds.

The two-slice induction setup: write a down-set on \([n+1]\) as
\[
        \mathcal D=\mathcal E\ \cup\ \{A\cup\{n+1\}:A\in\mathcal F\},
        \qquad \mathcal F\subseteq\mathcal E\subseteq2^{[n]}.
\]
Then \(Z(\mathcal D)=Z(\mathcal E)+\lambda Z(\mathcal F)\).  Membership
\(p{\bf1}_{n+1}\in P(\mathcal D)\) is equivalent to
\[
        p{\bf1}_n\in (1-p)P(\mathcal E)+pP(\mathcal F).
\]
A quantitative induction would follow from a lower bound on the largest
diagonal point \(T(\mathcal H)=\sup\{t:t{\bf1}\in P(\mathcal H)\}\) in terms of
\(Z_\lambda(\mathcal H)\).  Simple continuous affine lower bounds cannot be
sharp at the dictatorship because \(T=0\) at measure \(1-p\), while adding even
a tiny atom may make \(T\) jump.

\subsection*{Base-interval and cap-vector heuristics}

After multiplicative closure, the key interval can be taken to be
\([1/6,1/3]\).  Numerical searches with the nontrivial constraint
\(\max_i w_i<p\) suggest that for \(p\in[1/4,1/3]\) the worst hard examples are
near
\[
        (p,p,p,1-3p)
\]
with the first three coordinates just below \(p\).  In this configuration the
success event is essentially ``at least two of the four displayed coordinates
are selected'', giving
\[
        g_4(p)=\Pr[{\rm Bin}(4,p)\ge2]
        =6p^2(1-p)^2+4p^3(1-p)+p^4,
\]
and \(g_4(p)\ge p\) on \([1/4,1/3]\).  A possible route is therefore to prove
the cap-vector inequality
\[
        \max_i w_i<p,\quad p\in[1/4,1/3]
        \quad\Longrightarrow\quad
        \Pr\!\left[\sum_i w_i v_i\ge p\right]\ge g_4(p).
\]
This resembles a bounded-weight extremal principle, but no proof is yet in
hand.  Merging two small weights changes the crossing probability with opposite
signs on different residual-threshold intervals, so naive merging is not
monotone.

\subsection*{Deterministic fixed-size subset approach}

For rational \(p=a/b\), coloring into \(b\) colors and selecting a fixed
\(a\)-set of colors would be proved by the deterministic statement that every
\(b\)-tuple \(x_i\ge0\), \(\sum x_i=1\), has at least
\((a/b)\binom ba\) \(a\)-subsets with sum at least \(a/b\).  This is exactly the
Manickam--Miklos--Singhi type assertion for \(y_i=x_i-1/b\).  It holds when
\(a\mid b\) by Baranyai partitions (recovering the reciprocal coloring proof),
and the pair lemma proves the case \(a=2\), \(b\ge6\).  It is false at
\((a,b)=(3,10)\), so a proof for all rationals must either use randomness of the
color loads or add Bernoulli mixture information beyond one fixed subset size.

\subsection*{Potential use of comparison inequalities}

Pinelis-type left-tail domination bounds for sums of independent nonnegative
variables give upper bounds in terms of first and second moments, but the
available crude substitutions appear too weak near \(p=1/3\), where the desired
bad-event bound is \(1-p=2/3\).  For normalized variables
\(X_i=a_i\xi_i\), \(\sum a_i=L=1/p\), \(a_i<1\), \(\xi_i\sim{\rm Bernoulli}(1/L)\),
one has \(E\sum X_i=1\) and \(\sum E X_i^2=L^{-1}\sum a_i^2\).  The comparison
bounds are promising when this second-moment parameter is small, but near
cap-vector configurations they do not by themselves give the sharp constant.



\subsection*{Round 5 additional structural notes}

Finite sharpness is already present in a trivial way: if some weight \(w_0\ge p\)
and the sum of all other weights is \(<p\), then
\[
   \{\sum_i w_i v_i\ge p\}=\{v_0=1\}
\]
and the probability is exactly \(p\).  The pivot-plus-dust construction remains
important because it gives near equality while keeping the pivot below \(p\).

A useful rational coloring criterion: for \(p=a/b\), if every nonnegative
\(b\)-tuple \(x_1,\ldots,x_b\) with sum \(1\) has at least
\[
   {a\over b}\binom ba=\binom{b-1}{a-1}
\]
\(a\)-subsets \(J\) satisfying \(\sum_{j\in J}x_j\ge a/b\), then \(a/b\in
\mathcal G\).  Proof: randomly color original indices with \(b\) colors; for a
fixed \(a\)-set of colors, membership is iid Bernoulli\((a/b)\), and symmetry
converts the deterministic color-level count into the desired probability.
The cases \(a=1\) and \(a=2\) give the reciprocal points and the \(2/b\) family.
For general \(a\) this criterion is a bounded-below Manickam--Miklos--Singhi
type statement: with \(y_j=x_j-1/b\), the \(y_j\)'s sum to zero and satisfy
\(y_j\ge -1/b\), and one asks for many \(a\)-subsets of nonnegative \(y\)-sum.
The lower bound \(y_j\ge-1/b\) is stronger than the classical MMS hypotheses
and may be relevant, but a general proof for \(b\ge3a\) was not found.

The naive deterministic cyclic interval strengthening of the coloring criterion
is false.  For \(p=a/b=2/7\), take normalized color loads approximately
\[
   (1,0,1,0,1,0,1/2)
\]
with total \(7/2\) and threshold \(1\).  After a small downward perturbation of
the three \(1\)'s and adding the mass to the last coordinate, only one cyclic
interval of length \(2\) has load at least \(1\), although the desired count
would be \(2\).  Thus any cyclic proof must average or pair configurations; it
cannot rely on a pointwise interval count.

For the pivot functional
\[
 \Phi_a(Y)=\Pr[Y\ge p]-\lambda\Pr[Y<p-a],\qquad \lambda={p\over1-p},
\]
exposing a weight \(x\) and writing \(Y=Z+xv\) gives the exact identity
\[
\Phi_a(Z+xv)
=(1-p)\Bigl(A+\lambda B-\lambda C-\lambda^2D\Bigr),
\]
where
\[
 A=\Pr[Z\ge p],\quad
 B=\Pr[Z\ge p-x],\quad
 C=\Pr[Z<p-a],\quad
 D=\Pr[Z<p-a-x].
\]
Thus an induction on weights should control the four-tail quantity
\[
 A+\lambda B-\lambda C-\lambda^2D.
\]
The scalar pivot inequality for the residual only gives \(A\ge \lambda D\) and
does not control the possible deficit \(B-C\).

A related buffered candidate invariant is
\[
F_{a,h}(r)=
{\bf1}_{r\ge p}
-\lambda{\bf1}_{r<p-a}
+p{\bf1}_{p-h\le r<p}
+\lambda p{\bf1}_{p-a-h\le r<p-a}.
\]
If \(Y=R+h v\), then
\[
\Phi_a(Y)=\mathbb E F_{a,h}(R).
\]
This exactly records the two compensating boundary strips created when one
peels off a weight.  However, a direct pointwise induction
\((1-p)F_{a,h}(r)+pF_{a,h}(r+y)\ge F_{a,h+y}(r)\) is false for sample
parameters, so any use of this invariant must exploit distributional
constraints, not pointwise domination.

Complement/intersection viewpoint: with \(q=1-p\), the forbidden upper family
\[
  \mathcal U=\{A:\sum_{i\in A}w_i>q\}
\]
has complements of weight \(<p\).  Hence any \(k\) members of \(\mathcal U\)
intersect whenever \(kp\le1\); in particular for \(p\le1/3\) the family is
3-wise intersecting.  Frankl-type biased theorems for \(k\)-wise intersecting
families recover the reciprocal points when \(p=1/k\), but do not by themselves
cover non-reciprocal \(p\), because arbitrary \(k\)-wise intersecting families
can have too large \(q\)-biased measure once \(q>1-1/k\).  The quota/weighted
threshold structure is the missing extra ingredient.

Convex-hull formulation detail: for a down-set \(\mathcal D\),
\(P(\mathcal D)=\conv\{1_A:A\in\mathcal D\}\) is coordinatewise down-closed.
If \(x\in P(\mathcal D)\) is represented as marginals of a random
\(A\in\mathcal D\) and \(0\le y\le x\), thin each coordinate \(i\in A\) with
probability \(y_i/x_i\).  The thinned set remains in \(\mathcal D\) and has
marginals \(y\).









\section*{Round 4 additional working notes: pivot form and failed cap-vector route}

The intended classification remains
\[
        \mathcal G=[0,1/3]\cup\{1/2,1\},
\]
where \(\mathcal G\) denotes the set of \(p\) for which the inequality holds
for every finite probability vector of weights.  The counterexamples with
three equal weights for \(1/3<p<1/2\) and two equal weights for \(1/2<p<1\)
settle the negative direction.

\subsection*{Complement formulation}

For \(q=1-p\) and \(\eta_i=1-v_i\), the small-\(p\) lemma is equivalent to
\[
        q\ge 2/3,\qquad
        \Pr\!\left[\sum_i w_i\eta_i>q\right]\le q .
\]
Thus the missing assertion is a sharp strict upper-tail-at-the-mean bound for
convex combinations of iid Bernoulli\((q)\)'s with \(q\ge2/3\).  This is the
hypotenuse \(t=q\) in the Fokkink--Meester--Pelekis notation
\(\pi(q,t)=\sup_\gamma \Pr(S_\gamma\ge t)\), but with a strict inequality in
the event.  FMP prove several neighboring rectangles/isolated points; their
paper explicitly leaves triangular gaps along this hypotenuse.  The strict
version may be easier than the non-strict one because atoms at \(q\) are
ignored (for example \(q=2/3\), equal thirds have a large atom/tail at \(2/3\)
but small strict tail).

\subsection*{Pivot two-tail equivalence}

For \(0<p\le1/3\), the small-\(p\) lemma is equivalent to the following
two-tail inequality: for every \(0\le a\le p\), every finite family of weights
\(u_i\ge0\) with total \(1-a\), and iid Bernoulli\((p)\) variables,
\[
        \Pr\Big[\sum_i u_i v_i\ge p\Big]
        \ge {p\over 1-p}\,
        \Pr\Big[\sum_i u_i v_i<p-a\Big].        \tag{Piv(a)}
\]
Proof: \(a=0\) gives the original inequality by solving
\(U\ge\frac{p}{1-p}(1-U)\).  Conversely, assuming the original inequality,
add a new pivot weight \(a\) to the \(u_i\)'s.  Conditioning on the pivot gives
\[
p\Pr[Y\ge p-a]+(1-p)\Pr[Y\ge p]\ge p,
\]
which rearranges to \((\mathrm{Piv}(a))\).

This pivot form may be the best induction target.  If \(a=w_1\) is a largest
weight in an original problem and \(Y\) is the contribution of the remaining
weights, then the desired inequality is exactly
\[
(1-p)\Pr[Y\ge p]\ge p\Pr[Y<p-a].
\]
The thresholds are centered around \(E Y=p(1-a)\):
\[
p-EY=pa,\qquad EY-(p-a)=a(1-p),
\]
so the upper threshold is closer to the mean than the lower threshold by the
factor \(p/(1-p)\).

\subsection*{Sharpness and cap-vector warning}

The finite cap-vector lower bound
\[
\Pr\Big[\sum_i w_i v_i\ge p\Big]\ge
\Pr[\mathrm{Bin}(4,p)\ge2]\qquad (p\in[1/4,1/3])
\]
is false globally.  The actual near-equality mechanism is pivot plus dust.
Fix \(0<\varepsilon<p\), set \(a=p-\varepsilon\), and take
\[
        w_0=a,\qquad
        w_1=\cdots=w_N=(1-a)/N.
\]
With \(Y_N=(1-a)N^{-1}\mathrm{Bin}(N,p)\),
\[
\Pr[a v_0+Y_N\ge p]
 =
p\Pr[Y_N\ge\varepsilon]+(1-p)\Pr[Y_N\ge p].
\]
Since \(Y_N\to p(1-a)=p(1-p+\varepsilon)\) in probability and
\[
        \varepsilon<p(1-p+\varepsilon)<p
\]
for \(0<\varepsilon<p\), this probability tends to \(p\).  Therefore any proof
of the full interval must allow an infimum that is approached by one pivot and
an increasingly deterministic dust cloud; it cannot prove a uniform stronger
finite-support bound such as the four-cap binomial tail.

\subsection*{Local split/merge calculation for the pivot functional}

Let
\[
\Phi_a(w)=
\Pr\Big[\sum_i w_i v_i\ge p\Big]
-\lambda \Pr\Big[\sum_i w_i v_i<p-a\Big],
\qquad \lambda={p\over1-p},
\]
for weights of total \(1-a\).  The conjectured pivot inequality is
\(\Phi_a(w)\ge0\).  Replacing two weights \(x\ge y\) by a merged weight
\(c=x+y\), conditional on the rest being \(S=s\), the split-minus-merged
difference for the tail \(\{\cdot\ge T\}\) is \(p(1-p)h(T-s)\), where
\[
h(r)=
\begin{cases}
  1, & 0<r\le y,\\
  0, & y<r\le x,\\
 -1, & x<r\le x+y,\\
  0, & \text{otherwise}.
\end{cases}
\]
Thus the conditional split-minus-merged difference for \(\Phi_a\) is
\[
        p(1-p)\bigl(h(p-s)+\lambda h(p-a-s)\bigr).
\]
This is not pointwise nonnegative, and numerical checks show that arbitrary
pairwise merging is not monotone for \(\Phi_a\).  However, random tests did not
find any violation of \(\Phi_a\ge0\) for \(p\le1/3\).

\subsection*{Continuous rotation/orbit attempt}

Represent Bernoulli\((p)\) as \(1_{U_i\in[0,p]}\) with \(U_i\) uniform on the
circle.  For fixed starting points, consider the orbit under common rotation
and the weighted coverage function
\[
        f(t)=\sum_i w_i 1_{\{U_i+t\in[0,p]\}}.
\]
If every orbit satisfied
\[
        |\{t:f(t)\ge p\}|\ge p,
\]
then the theorem would follow by averaging.  This deterministic orbit
statement is true for the reciprocal cyclic proof but is false in general:
random searches found configurations of four arcs of length \(p=0.3\) and
weights near \(0.25\) for which the measure of \(\{f\ge p\}\) is about
\(0.268<p\).  Hence an orbitwise proof by one-dimensional rotations cannot
work without an additional averaging or higher-dimensional device.

\subsection*{Computational sanity checks}

Brute-force enumeration for random weights of size up to \(8\) and many random
parameters \(a\in[0,p]\) found no violation of \((\mathrm{Piv}(a))\) for
\(p=0.2,0.3,1/3\).  In particular, the one-weight case has equality:
if the single weight is \(1-a\ge p\), then
\[
\Pr[Y\ge p]=p,\qquad \Pr[Y<p-a]=1-p,
\]
so \(\Phi_a=0\).  This is consistent with pivot-plus-dust as the limiting
extremal structure.

Potential next proof strategies:
\begin{itemize}
\item Prove \((\mathrm{Piv}(a))\) by induction with a stronger two-threshold
invariant.  A naive scalar Bellman bound fails because adaptive residual
strategies can exceed the desired bound; the residual must be controlled
simultaneously at the two thresholds \(p-a\) and \(p\).
\item Work in the complement formulation and prove that every weighted
threshold intersecting family
\(\{A:\sum_{i\in A}w_i>q\}\), \(q\ge2/3\), has \(q\)-biased measure at most
\(q\).  This is false for arbitrary intersecting families when \(q>1/2\), so
the linear-threshold/quota structure is essential.
\item Search for a known sharp theorem on weighted Bernoulli sums exceeding
their mean for \(q\ge2/3\).  The closest located literature is
Fokkink--Meester--Pelekis on optimizing stakes; it does not appear to close the
strict hypotenuse case directly.
\end{itemize}


\subsection*{Round 6 explorations: weighted-threshold/intersection viewpoint}

For the small-\(p\) lemma it is useful to normalize
\[
  a_i=w_i/p,\qquad L=1/p,\qquad \xi_i\sim{\rm Bernoulli}(1/L).
\]
The nontrivial case \(w_i<p\) becomes \(0\le a_i<1\), \(\sum_i a_i=L\ge3\), and
\[
  \Pr\!\left[\sum_i a_i\xi_i\ge1\right]\ge\frac1L .
\]
By multiplicative closure it remains enough to prove \(L\in[3,6]\)
(\(p\in[1/6,1/3]\)).

Complement/down-set structure: for
\[
  \mathcal D_w(p)=\{A:\sum_{i\in A}w_i<p\}
\]
and \(r=\lfloor1/p\rfloor\), no \(r\) sets in \(\mathcal D_w(p)\) cover the
ground set, since the union has weight at most the sum of the individual
weights, which is \(<rp\le1\).  The complement family
\[
  \mathcal H_w(p)=\{B:\sum_{i\in B}w_i>1-p\}
\]
is therefore \(r\)-wise intersecting and is a weighted-threshold family.  The
small-\(p\) lemma is equivalent to
\[
  \mu_{1-p}(\mathcal H_w(p))\le1-p
\]
for these weighted-threshold \(r\)-wise intersecting families.  A plain
``\(r\)-wise intersecting'' theorem is not sufficient: e.g. for \(r=3\), the
family of all subsets of size \(>2n/3\) is 3-wise intersecting and has
\(\mu_q\) close to \(1\) when \(q>2/3\).  The additional fact that the threshold
level is \(q=1-p\), with the same weights that certify the intersection
property, is essential.

Connection with Fokkink--Meester--Pelekis: the complement form is a
``favorable odds'' upper-tail problem on the hypotenuse \(t=q\), namely
\[
  \Pr\!\left[\sum_i w_i\eta_i>q\right]\le q,\qquad q=1-p\ge2/3.
\]
FMP prove neighboring rectangles and the special points \(q=(k+1)/(k+2)\), but
their published theorem does not by itself fill the whole hypotenuse
\(q\in[2/3,1]\).  Thus a direct proof here would amount to a special
hypotenuse theorem for weighted threshold families.

A potentially useful rational/cyclic representation: for \(p=c/b\), assign each
weight independently to a random point of \(\mathbb Z/b\).  For each start \(j\),
let \(Y_j\) be the load in the length-\(c\) cyclic interval.  Then each \(Y_j\)
has the target distribution, and the desired inequality is
\[
  \mathbb E\,|\{j:Y_j\ge c/b\}|\ge c
\]
(or, in complement form, the expected number of length-\(b-c\) intervals with
load \(> (b-c)/b\) is at most \(b-c\)).  The corresponding deterministic
statement for every load vector is false; any proof must use randomness of the
assignment.  This is analogous to the injection/averaging arguments in FMP's
proofs at special rational points.

Pinelis left-tail domination was checked as a possible black box.  The coarse
version depending only on \(m=\mathbb E S\) and \(s=\sum\mathbb E X_i^2\)
compares \(S\) with \(s/m\) times a Poisson random variable.  At threshold
\(m=p\) this gives \(P({\rm Pois}(\lambda)<\lambda)\), which is too weak near
one large weight (for \(\lambda=p\), it gives \(e^{-p}>1-p\)).  The
individual-moment version appears essentially sharp for the original variables
and does not immediately improve the problem.

MathOverflow/literature note: a related question on lower bounds for
\(\Pr[\sum a_i\xi_i\ge p\sum a_i]\) has answers giving absolute-constant
bounds using Paley--Zygmund, Berry--Esseen, and FKG decompositions.  These are
not sharp enough for the exact \(p\) bound but may be useful for regimes where
all weights have high multiplicity or very small maximum weight.



\subsection*{Additional structural reductions}

\paragraph{Closedness of the good set.}
Let \(\mathcal G\) be the set of \(p\)'s for which the weighted Bernoulli
inequality holds universally.  This set is topologically closed.  For fixed
finite weights \(w\), the moving-threshold tail
\[
  F_w(r)=\mu_r\{A:w(A)\ge r\}
\]
has only finitely many threshold changes.  If \(p_n\to p\) with \(p_n\in
\mathcal G\), then along a subsequence \(p_n\le p\) one has
\(F_w(p_n)\to \mu_p\{w(A)\ge p\}=F_w(p)\), while along a subsequence
\(p_n>p\) one has \(F_w(p_n)\to\mu_p\{w(A)>p\}\le F_w(p)\).  In either case
\(F_w(p)\ge p\).  Thus a dense subset of positive values in \([0,1/3]\) would
settle the small-\(p\) lemma.

\paragraph{Unit-gap form and its exact relation to the pivot.}
For \(q=1-p\), the small-\(p\) lemma is equivalent to the following form:
for all \(0\le b_i\le1\), \(B=\sum b_i\), and iid
\(\zeta_i\sim{\rm Bernoulli}(p)\),
\[
 \Pr\!\left[\sum_i b_i\zeta_i\ge p(B+1)\right]
 \ge {p\over q}\,
 \Pr\!\left[\sum_i b_i\zeta_i<p(B+1)-1\right]. \tag{UG}
\]
Indeed, (UG) is obtained by applying the original inequality to the coefficient
vector \((1,b_1,\ldots,b_n)\) and conditioning on the unit coefficient.  In the
other direction, if the original normalized weights all have maximum
\(a<p\), divide the non-pivot weights by \(a\).  Then \(a(B+1)=1\), and (UG)
becomes the two-tail pivot inequality
\[
 \Pr[Y\ge p]\ge {p\over q}\Pr[Y<p-a].
\]
Conditioning on the pivot gives the original inequality.  Thus (UG) is not a
strictly weaker lemma; it is the original problem in the normalization where
the largest coefficient has been scaled to \(1\).  It may nevertheless be the
right Bellman/induction target because the two thresholds differ by exactly one
and every residual jump has size at most one.

\paragraph{Unit coefficients in the unit-gap form.}
When all \(b_i\in\{0,1\}\), say \(B=N\), (UG) reduces to a binomial tail ratio.
Let \(X\sim{\rm Bin}(N,p)\), \(Y=X+\zeta_0\sim{\rm Bin}(N+1,p)\), and
\(k=\lceil p(N+1)\rceil\).  Then
\[
 \Pr[Y\ge p(N+1)]
 =p\,\Pr[X\ge k-1]+q\,\Pr[X\ge k]
 =p+q\,\Pr[X\ge k]-p\,\Pr[X\le k-2].
\]
Thus the standard binomial inequality
\(\Pr[Y\ge (N+1)p]\ge p\) is exactly (UG) in this case.  This observation
explains why the reciprocal/equal-weight cases fit the same pivot framework.

\paragraph{Possible induction state.}
For scaled coefficients \(0\le c_i\le1\) with total \(C\) and at least one
coefficient equal to \(1\), the desired statement is
\[
 \Pr\left[\sum_i c_i\zeta_i\ge pC\right]\ge p.
\]
If \(C\le1/p\) the unit coefficient alone proves the claim.  The hard case is
\(C>1/p\), and conditioning on a unit coefficient leaves thresholds
\(pC-1=p(C-1)-q\) and \(pC=p(C-1)+p\) around the residual mean.  This is exactly
the asymmetric two-tail comparison in (UG).  A successful induction likely
needs to control the residual distribution simultaneously at these two
thresholds; one-parameter induction at the mean is insufficient.




\subsection*{Coalescence formulation for the complement problem}

For \(q=1-p\), the small-\(p\) lemma is equivalent to
\[
 \Phi_q(w):=\Pr_q\left[\sum_i w_i\eta_i>q\right]\le q,\qquad q\ge2/3.
\]
If two weights \(x\le y\) are coalesced to \(x+y\), and \(Z\) is the sum of
the remaining coordinates, then conditioning on \(Z\) gives the exact change
\[
 \Delta
 =\Phi_q(\ldots,x+y,\ldots)-\Phi_q(\ldots,x,y,\ldots)
 =pq\left(\Pr[q-x-y<Z\le q-y]-\Pr[q-x<Z\le q]\right).
\]
In failure variables with failure probability \(p\), the same formula is
\[
 \Delta=pq\left(\Pr[p-x\le Z<p]-\Pr[p-x-y\le Z<p-y]\right),
\]
with endpoint conventions adjusted to the strict event \(w(F)<p\).

The pairwise coalescence assertion ``some pair always has \(\Delta\ge0\)''
is false even when all weights are below \(p\).  Infinite equal-weight trap:
if \(n\ge4\), \(1/n<p<1/(n-1)\), \(q=1-p\), and
\(w_1=\cdots=w_n=1/n\), then \(w_i<p\) and
\[
 \Phi_q(w)=q^n+npq^{n-1}
\]
(empty or singleton failure set).  After coalescing any pair, the merged weight
\(2/n\) already exceeds \(p\), so allowed failure sets are empty or one of the
\(n-2\) unmerged singletons:
\[
 \Phi_q(w^{ij})=q^{n-1}+(n-2)pq^{n-2}.
\]
Thus
\[
 \Delta=pq^{n-2}\bigl(n-2-(n-1)q\bigr)<0.
\]
Concrete instance: \(p=.3,q=.7,w=(1/4)^4\), where
\(\Phi_q=0.6517\) and every pair-merge gives \(0.637\).

Possible salvage: prove a dichotomy saying either some pair has nonnegative
gain or the vector lies in a reciprocal-chamber configuration that can be
bounded directly by a binomial tail.  Block coalescence can also work where
pair coalescence fails: in the \(n=4,p=.3\) example, merging three of the four
weights gives a coordinate \(3/4>q\), hence value \(q\), above the original
\(0.6517\).  A useful exact block-gain formula: for a block \(B\) of total mass
\(c\), with internal random sum \(T_B\) and external sum \(Z\),
\[
 \Delta_B=\mathbb E_Z\left[
 q\,{\bf 1}_{\{Z+c>q\}}-\Pr(T_B>q-Z)
 \right].
\]
No general nonnegative proper-block theorem is proved.

\subsection*{Fokkink--Meester--Pelekis / bold-play formulation}

The complement inequality is a case of the FMP optimization problem
\[
 \pi(q,t)=\sup_{\sum c_i=1}\Pr_q\left[\sum_i c_i\beta_i\ge t\right].
\]
Our target is \(\Pr_q(\sum c_i\beta_i>q)\le q\).  For a fixed finite vector,
strict \(>q\) can be written as \(\ge t_*(c)\) where \(t_*(c)\) is the smallest
subset sum exceeding \(q\), so a theorem \(\pi(q,t)\le q\) for every \(t>q\)
would imply the small-\(p\) lemma.

FMP Theorem 4.4 proves bold play in rectangles
\[
 \frac{k}{k+1}<q\le \frac{k+1}{k+2}<t.
\]
This covers cases where the first attainable threshold above \(q\) jumps beyond
the next reciprocal \((k+1)/(k+2)\).  It leaves the near-hypotenuse strip
\[
 q<t\le \frac{k+1}{k+2},\qquad \frac{k}{k+1}<q<\frac{k+1}{k+2}.
\]
FMP Proposition 4.5 proves the boundary points
\(q=t=(k+1)/(k+2)\) for \(k>1\) (and identifies the exceptional
\(q=t\in[1/2,2/3]\) behavior).  For the original problem these boundary points
correspond to the reciprocal values \(p=1/(k+2)\).  The desired interval
\(p\le1/3\) is precisely the collection of upper-right hypotenuse strips
\(q\ge2/3\) that FMP note are suggested numerically but not fully settled in
their paper.

\subsection*{Bellman/unit-gap induction attempt}

The unit-gap inequality can be centered as follows.  For \(0\le b_i\le1\),
\(B=\sum b_i\), and \(\zeta_i\sim\Ber(p)\), put
\[
 S=\sum_i b_i(\zeta_i-p).
\]
Then (UG) is
\[
 \Pr[S\ge p]\ge \frac pq \Pr[S<-q].
\]
Thus one needs a two-tail comparison for a mean-zero sum whose positive jumps
are \(qb_i\), negative jumps are \(-pb_i\), and jump ranges are at most \(1\).

A tempting stronger induction would assert, for all shifts \(x\),
\[
 \Pr[S\ge x+p]\ge \frac pq\Pr[S<x-q].
\]
Conditioning on a coefficient \(b\) closes this family because it asks for the
same inequality for the residual sum at shifts \(x-qb\) and \(x+pb\).  However,
the zero-variable base case fails for \(x>q\).  The restricted domain
\(x\le q-pB\) is stable under conditioning but does not contain the required
point \(x=0\) when \(B>q/p\), which is exactly the hard large-dust regime.
This explains why a one-parameter induction fails.  A successful Bellman
argument likely needs an additional middle-interval term or a value function
depending on both the total remaining mass and the shift.

\subsection*{Influence/threshold-flow observation}

For a fixed weighted threshold down-set
\(D_t=\{A:w(A)<t\}\), define \(g(r)=\mu_r(D_t)\).  Russo's formula gives
\[
 -g'(r)=I_r(D_t),
\]
the total \(r\)-biased influence of the threshold.  The desired inequality is
\(g(t)\le1-t\) for \(t\le1/3\).  It would follow from
\(\int_0^t I_r(D_t)\,dr\ge t\).  The pointwise bound
\(I_r(D_t)\ge1\) is false for small \(r\) even when \(t\le1/3\): equal small
weights have very small initial influence.  Numerically, however,
\(I_t(D_t)\ge1\) appears to hold, with equality in dictator-type cases.  This
may be a useful endpoint/stability diagnostic but is not currently a proof.

\subsection*{Circle/interval representation}

In the normalized form \(L=1/p\), \(a_i=w_i/p\), \(\sum_i a_i=L\),
\(a_i\le1\), the desired statement is
\[
 \Pr\left[\sum_i a_i\xi_i\ge1\right]\ge1/L,\qquad \xi_i\sim\Ber(1/L).
\]
Equivalently, place independent uniform points \(U_i\) on a circle of
circumference \(L\); the load of a fixed unit interval has this distribution.
For a deterministic configuration the set of unit intervals with load at least
\(1\) need not have length at least \(1\) when \(L\) is nonintegral.  Numerical
counterconfigurations occur already for \(L=3.1\) and \(L=3.5\).  Thus the
integral reciprocal coloring proof does not extend by a pointwise continuous
circle lemma; any circle proof must use the randomness of the placements or a
more subtle averaging.

\subsection*{Disjoint-copy coupling attempt}

For \(p\le1/3\), one can couple three Bernoulli\((p)\) selections disjointly
coordinatewise by assigning each coordinate to color \(1,2,3\) with probability
\(p\) each and to a leftover color with probability \(1-3p\).  The three color
loads have the desired marginal distribution.  A deterministic implication from
the total active load to the number of color classes of load at least \(p\) is
too weak: three loads can all be below \(p\) while their total is arbitrarily
close to \(3p\).  Bounds based only on the active total therefore do not recover
the desired \(p\) lower bound.



\subsection*{Round 10 additional structural observations}

\subsubsection*{A deterministic coloring strengthening fails}

A natural strengthening of the coloring criterion would be: for all
\(a/b\le 1/3\) and all nonnegative \(x_1,\ldots,x_b\) with sum \(1\), at
least \(\binom{b-1}{a-1}\) \(a\)-subsets have sum at least \(a/b\).  This is
false.  Put \(b=3a+1\), take three zero coordinates and \(3a-2\) coordinates
equal to \(1/(3a-2)\).  Then an \(a\)-subset has sum at least
\(a/(3a+1)\) iff it avoids all three zero coordinates, so the number of good
sets is
\[
        \binom{3a-2}{a},
\]
which is smaller than \(\binom{3a}{a-1}\) for \(a\ge3\).  The first case is
\(a=3,b=10\), where the count is \(35<36\).  Thus the all-\(a\)-subsets
averaging/coloring criterion cannot prove the whole interval \([0,1/3]\).
This is the bounded-below version of the familiar \(n=3k+1\) obstruction to
the naive Manickam--Miklos--Singhi strengthening.

The original Bernoulli inequality still holds comfortably for this vector:
for \(p=a/(3a+1)\) it reduces to
\[
   \Pr\{\operatorname{Bin}(3a-2,p)\ge a\}\ge p,
\]
with a margin already positive at \(a=3\).

\subsubsection*{Convex-hull/partition-function target}

The clean abstract target remains the following.  For
\(\lambda=p/(1-p)\le1/2\), if \(\mathcal D\subseteq2^{[n]}\) is a down-set and
\(p{\bf1}\notin P(\mathcal D)=\operatorname{conv}\{1_A:A\in\mathcal D\}\),
then
\[
       Z_\lambda(\mathcal D):=\sum_{A\in\mathcal D}\lambda^{|A|}
       \le (1+\lambda)^{n-1}.
\]
Equivalently, every down-set with \(Z_\lambda(\mathcal D)>(1+\lambda)^{n-1}\)
should admit a distribution on its members with all coordinate marginals at
least \(p\) (then thinning gives exactly \(p\)).  The extremal example is the
dictatorship \(\mathcal D=2^{[n]\setminus\{i\}}\), which has equality and
misses \(p{\bf1}\).  Brute force over all down-sets for \(n\le4\) at
\(p=1/3\) gives only dictatorships as maximizers among down-sets missing
\(p{\bf1}\).

The two-section induction formulation is: write
\[
 \mathcal D=\mathcal E\cup\{A\cup\{n\}:A\in\mathcal F\},\qquad
 \mathcal F\subseteq\mathcal E.
\]
Then \(Z(\mathcal D)=Z(\mathcal E)+\lambda Z(\mathcal F)\), and
\[
 p{\bf1}_{n}\in P(\mathcal D)
 \Longleftrightarrow
 p{\bf1}_{n-1}\in (1-p)P(\mathcal E)+pP(\mathcal F).
\]
A sufficient induction lemma would therefore be:
if the displayed Minkowski sum misses \(p{\bf1}_{n-1}\), then
\(Z(\mathcal E)+\lambda Z(\mathcal F)\le(1+\lambda)^{n-1}\).  The simple
one-section induction only shows that if \(p{\bf1}\notin P(\mathcal F)\) then
\(Z(\mathcal F)\le(1+\lambda)^{n-2}\); it does not control the possible
large value of \(Z(\mathcal E)\).

\subsubsection*{Intersection-theoretic endpoint information}

For \(p_{\rm orig}\le1/3\), \(q=1-p_{\rm orig}\ge2/3\), and
\[
  \mathcal H_w=\{A:w(A)>q\},
\]
the family \(\mathcal H_w\) is \(r\)-wise intersecting for
\(r=\lfloor1/p_{\rm orig}\rfloor\).  General \(r\)-wise intersection theorems
are not by themselves enough in the open strips.  The known biased-measure
star bound for \(r\)-wise intersecting families applies up to bias
\((r-1)/r\), whereas here \(q\) can be as large as \(r/(r+1)\) while the
family is not necessarily \((r+1)\)-wise intersecting.  At the endpoint
\(p_{\rm orig}=1/3\) this recovers the familiar \(3\)-wise-intersecting
measure bound at bias \(2/3\), but the reciprocal point was already covered by
the coloring argument.

Equal-weight threshold families illustrate why a pure \(r\)-wise-intersection
argument is too weak: for \(q>2/3\), the family
\(\{|A|>qn\}\) is \(3\)-wise intersecting by the union bound, and its
\(\mu_q\)-measure is about \(1/2\), below \(q\); but arbitrary
\(3\)-wise-intersecting families can have much larger \(\mu_q\)-measure once
\(q>2/3\).  The missing ingredient is the halfspace/weighted-threshold
structure at the same bias as the threshold.

\subsubsection*{Exponential-clock representation}

Writing \(v_i(t)=1_{\{E_i\le t\}}\) with iid exponential clocks and
\(p=1-e^{-t}\), the original probability is
\[
   \Pr\left[\sum_i w_i v_i(t)\ge p\right].
\]
For a fixed random permutation of the clock order, let \(K_p\) be the first
rank at which the cumulative weight in that order reaches \(p\).  The order
statistics are independent of the permutation, and the Bernoulli subset at
time \(t\) can equivalently be obtained by taking the first \(N\) elements of
the permutation, where \(N\sim\operatorname{Bin}(m,p)\) is independent of the
order.  Hence
\[
   \Pr\left[\sum_i w_i v_i\ge p\right]
   =\mathbb E\big[\Pr\{\operatorname{Bin}(m,p)\ge K_p\mid K_p\}\big].
\]
A proof would follow from a sufficiently sharp universal constraint on the
distribution of the random crossing rank \(K_p\).  Simple constraints such as
bounds on \(\mathbb E K_p\) appear too weak, but this representation may be a
useful ballot/cyclic-shift viewpoint.


\end{document}
